%----------------------------------------------------------------------------------------
%	RESEARCH
%----------------------------------------------------------------------------------------

\section{Research Experience}

\researchentry
  {Starspot Inference with Light Curve Inversion Techniques}
  {American Museum of Natural History}
  {Dr. Lucy Lu \enspace\textbar\enspace \textit{M.S. Thesis Researcher}}
  {Aug 2023 -- Sep 2025}
  {\thesislink{https://academicworks.cuny.edu/gc_etds/6480/}}
\begin{itemize}[
  leftmargin=1em,
  label={--},
  itemsep=3pt,
  topsep=0pt,
  parsep=0pt,
  partopsep=0pt
]
  \item Analyzed 1,000 synthetic, evolving starspot light curves generated
  with \texttt{butterpy} and implemented forward- and inverse-modeling
  workflows with \texttt{fleck} and \texttt{starry}.

  \item Evaluated recovery of stellar rotation, activity indicators, and
  surface structure across variations in photometric noise, viewing
  inclination, and spot configuration.

  \item Demonstrated that rotation and statistical activity indicators could
  be recovered reliably, while individual starspot locations and surface
  maps remained fundamentally degenerate.
\end{itemize}

\vspace{0.4em}

\researchentry
  {Towards Mapping Brown Dwarf and Giant Exoplanet Atmospheres}
  {American Museum of Natural History}
  {Dr. Johanna Vos \enspace\textbar\enspace \textit{BDNYC Researcher}}
  {May 2021 -- Aug 2023}
  {\githubrepo{https://github.com/MohammadRefat23/StarryAtmospheres}}

\begin{itemize}[
  leftmargin=1em,
  label={--},
  itemsep=3pt,
  topsep=0pt,
  parsep=0pt,
  partopsep=0pt
]

\item Modeled three-dimensional atmospheric simulations and 20-hour
Spitzer light curves of two isolated planetary-mass objects using
\texttt{starry} spherical-harmonic surface maps.

\item Compared surface maps with spherical-harmonic degrees $l=1$, 5, and
10, identifying $l=5$ as sufficient to reproduce the relevant light-curve
structure without unnecessary map complexity.

\item Recovered the overall light-curve morphology of 2M0030 and showed that
time evolution limited a single static-map fit to 2M0642, motivating
rotation-by-rotation atmospheric mapping.

\end{itemize}

\vspace{0.4em}

\researchentry
  {Chemodynamically Characterizing the Jhelum Stellar Stream with APOGEE-2}
  {Sloan Digital Sky Survey Faculty and Student Team}
  {Dr. Allyson Sheffield \enspace\textbar\enspace \textit{SDSS FAST Student Researcher}}
  {May 2020 -- May 2021}
  {%
    \githubrepo{https://github.com/MohammadRefat23/jhelum-stellar-stream}
    \hspace{0.15em}
    \paperlink{https://doi.org/10.3847/1538-4357/abee93}
  }

\begin{itemize}[
  leftmargin=1em,
  label={--},
  itemsep=3pt,
  topsep=0pt,
  parsep=0pt,
  partopsep=0pt
]
\item Analyzed 1,495 APOGEE-2 targets, selecting 289 red giants and combining
Gaia astrometry, stellar kinematics, and chemical abundances to identify
eight candidate Jhelum stream members.

\item Identified one new metal-poor red-giant-branch member
($[\mathrm{Fe}/\mathrm{H}]=-2.2$) and compared its abundance and orbital
patterns with candidate Gaia--Enceladus--Sausage stars.

\item Coauthored a peer-reviewed article in \textit{The Astrophysical Journal}
finding that the member's abundances were consistent with an accreted
dwarf-galaxy progenitor, while not excluding a globular-cluster origin.
\end{itemize}

\vspace{0.4em}

\researchentry
  {Looking at Star Formation Through Chemistry}
  {American Museum of Natural History}
  {Dr. Allyson Sheffield \enspace\textbar\enspace \textit{AstroCom NYC Scholar}}
  {Aug 2018 -- May 2020}
  {}

\begin{itemize}[
  leftmargin=1em,
  label={--},
  itemsep=3pt,
  topsep=0pt,
  parsep=0pt,
  partopsep=0pt
]
\item Processed and normalized high-resolution echelle spectra with
\texttt{SMHr}, measured absorption-line equivalent widths, and iteratively
fit stellar atmospheric parameters and elemental abundances.

\item Validated the workflow on Arcturus, recovering
$T_{\mathrm{eff}}=4300\,\mathrm{K}$, $\log g=1.6$, and
$[\mathrm{Fe}/\mathrm{H}]=-0.65$---within 50 K and 0.07 dex of published
values.

\item Established a proof-of-concept chemical-tagging workflow for
distinguishing in-situ, kicked-out, and accreted formation scenarios for
Galactic red-giant populations.
\end{itemize}

\vspace{0.4em}

\researchentry
  {Two-Point Statistics in the Star-Forming Interstellar Medium}
  {Flatiron Institute, Center for Computational Astrophysics}
  {Dr. Chang-Goo Kim \enspace\textbar\enspace \textit{AstroCom NYC Scholar}}
  {Jun 2018 -- Aug 2018}
  {\githubrepo{https://github.com/MohammadRefat23/tigress-metallicity-correlations}}

\begin{itemize}[
  leftmargin=1em,
  label={--},
  itemsep=3pt,
  topsep=0pt,
  parsep=0pt,
  partopsep=0pt
]
\item Developed Python routines to compute and visualize two-point
metallicity correlations in TIGRESS magnetohydrodynamic simulations of the
star-forming interstellar medium.

\item Measured a characteristic half-correlation length of
$120\pm25\,\mathrm{pc}$ and estimated a turbulent metal-mixing timescale of
$\sim10\,\mathrm{Myr}$.

\item Found no direct relationship between the measured correlation length
and either the supernova rate or gas velocity dispersion, constraining the
drivers of chemical inhomogeneity in the simulated gas.
\end{itemize}
